\documentclass[11pt]{article}

\usepackage[review]{acl}

\usepackage{times}
\usepackage{latexsym}
\usepackage[T1]{fontenc}
\usepackage[utf8]{inputenc}
\usepackage{microtype}
\usepackage{inconsolata}
\usepackage{graphicx}
\usepackage{booktabs}

\title{Non-Exchangeability in KV-Cache Relay: From Fixed Pipelines to Recursive, Model-Driven Delegation}

\author{Anonymous ACL submission}

\begin{document}
\maketitle

\begin{abstract}
Multi-agent LLM systems that relay key-value (KV) cache between agents
usually assume that some part of that cache -- which layer, which
channel, which question's content -- can be swapped out or reused. This
assumption is rarely tested directly. We test it with a causal audit: we
replace the relayed cache with zeroed content, random content, or a
different question's real content, and check whether accuracy tracks the
change. We run this audit across three structurally different
multi-agent topologies: a fixed sequential three-agent pipeline, a static
fan-in decomposition, and -- our central contribution -- a
\emph{dynamic, model-driven} delegation loop in the style of Recursive
Language Models (RLMs) \citep{rlm2026}. In this last topology, the model
itself writes and runs code to decide what to delegate and when, and a
sub-call's answer is relayed as real KV cache instead of RLM's own
decoded-text channel. The same three-tier ordering holds on every
topology we test: real content beats wrong-but-real content, which beats
no real content. This holds with statistical significance everywhere,
including the model-driven topology -- to our knowledge the first causal
audit of content-specific KV relay under model-driven, rather than
predetermined, delegation. The fan-in topology gives the cleanest result
in our study ($p=0.0001$ for real vs.\ wrong-but-real content, the
strongest instance of this comparison we observe). On the recursive
topology, we find a second, complementary signal: the model's own
\emph{behavior} -- how often it commits to an answer within its turn
budget, not just its final answer -- tracks the same ordering. This
behavioral signal resolves one comparison (wrong-but-real vs.\ no real
content) that final-answer accuracy alone cannot. We ground this
progression in a fixed-pipeline foundation with two headline results of
its own: relayed KV carries real, content-specific information, not just
a non-empty cache (established on the simplest, most rigorously tested
topology, then shown to generalize); and a calibration procedure's own
confidence does not predict -- and actively misleads on -- which of two
model architectures is safe to compress along the depth axis. We report
where four published KV-relay efficiency techniques transfer to this
setting and where they do not. We also report honestly where our
recursive-delegation results are less mature than our fixed-pipeline
evidence: a real non-determinism bug in the recursive prototype's
decoding loop, which we diagnosed and fixed, confirmed reproducible under
a forced-determinism rerun; and an accuracy-over-text-relay comparison
that remains statistically unresolved, even though the causal
content-dependence claim is not.
\end{abstract}

\section{Introduction}

Relaying key-value cache between agents in a multi-agent LLM system rests
on an assumption that is rarely stated and, to our knowledge, never
directly tested: that some dimension of the cache -- which transformer
layer, which channel within a head, which question's content -- can be
treated as generic or interchangeable for the purpose of compression,
reconstruction, or reuse. Prior audits of this assumption, ours included
in earlier form, fix the multi-agent structure itself: a predetermined
sequence or decomposition of who talks to whom, decided in advance by the
system's designer, not by the model. Recursive Language Models
\citep{rlm2026} describe a materially different regime already seeing
real use: a root model writes and executes code in a live Python REPL to
decide, for itself, what to delegate to a sub-call and when, rather than
following a fixed script. RLM's own sub-call return channel is
text -- a decoded string the root model may or may not even read. The
natural question this paper asks is whether the same non-exchangeability
result holds when that return channel is real KV cache instead, and,
more fundamentally, whether \emph{any} causal content-dependence claim
about relayed KV survives the move from a designer-fixed topology to one
the model itself constructs on the fly. No existing audit of KV/latent
reuse -- ours or the published work we compare against in
Section~\ref{sec:related} -- has tested this.

This paper answers that question by building the claim up through three
structurally distinct topologies, each testing more of the design space
than the last. First, within a fixed sequential three-agent
(Reasoner--Verifier--Finalizer) pipeline, we test three structurally
distinct axes of the exchangeability assumption and find it false along
all three: dropped transformer layers are not safely approximated by
their neighbors; key vectors' channels cannot be freely mixed once
rotary position embeddings have imposed position-dependent structure on
them; and a receiving agent's accuracy depends on the specific content of
the relayed cache, not merely on receiving \emph{some} non-empty cache,
established with direct causal evidence rather than inferred from
downstream accuracy alone. Second, we show the content-identity result
generalizes to a static fan-in decomposition, where an aggregator
receives two parallel sources rather than one sequential handoff --
still designer-fixed, but structurally distinct, and here the causal
ordering is the cleanest and most decisive we observe anywhere in this
paper. Third, and central to this paper's contribution, we show it
generalizes further still to a dynamic, model-driven delegation loop
built in RLM's own style, with KV cache substituted for RLM's
decoded-text return channel: the model itself decides, turn by turn,
whether and what to delegate, and we causally audit the KV it receives
back exactly as we do in the two fixed topologies. The same three-tier
ordering (real content $>$ wrong-but-real content $>$ no real content)
holds there too, statistically significantly, including one comparison
that final-answer accuracy alone cannot yet resolve but a complementary
behavioral signal -- how often the model commits to an answer within its
turn budget, not only what it eventually answers -- resolves cleanly.

We call the pattern underlying all three topologies
\textbf{non-exchangeability}: efficiency and delegation techniques for
KV-cache relay fail, or lose their causal grounding, in proportion to how
much they implicitly assume exchangeability along an axis the underlying
representation does not actually support. Within the fixed-pipeline
foundation specifically, the sharpest single result is that this failure
is invisible from the inside: a calibration procedure's own confidence
does not predict which of two model architectures is safe to compress
along the depth axis, and fails in the \emph{wrong} direction -- the
architecture that collapses harder under layer-selection produces the
\emph{more} confident-looking signal (Section~\ref{sec:results}, Finding
6). We demonstrate the broader pattern by diagnosing and, where possible,
correcting four published KV-relay efficiency techniques within our
fixed pipeline: a cross-agent offset-correction method \citep{kvcomm2025}
fails because it assumes a delta observed for one question transfers to
another; a rotation-based quantization scheme
\citep{turboquant2026,polarquant2026} fails specifically on key vectors
because it assumes head dimensions are interchangeable, disrupting RoPE's
position-dependent structure; a per-channel quantization fix
\citep{kivi2024} resolves this by respecting that structure instead; and
layer-selection's cost is architecture-dependent in a way its own
calibration signal does not predict. We report honestly where our
newest evidence is not yet as mature as this foundation: the recursive
prototype's decoding loop initially showed a real, diagnosed
non-determinism absent from the fixed pipeline -- since fixed and
confirmed reproducible under a forced-determinism rerun (Results) -- and
its accuracy-over-text-relay comparison remains statistically unresolved
even where its causal content-dependence claim is not.

Three recent papers anticipate pieces of the fixed-pipeline argument from
different angles -- that compression misses task-specific information
\citep{pitfallskv2026}, that layer redundancy depends on the calibration
objective rather than being a fixed property \citep{rethinkredundancy2026},
and that adaptive, non-fixed-ranking methods can generalize across
architectures where fixed ones might not \citep{adak2026}. A fourth,
very recent line of work \citep{canonicalmerge2026} builds real fan-in KV
merging on the same models and a closely related task to our own fan-in
topology, but tests only downstream accuracy, not causal
content-dependence, and does not consider model-driven delegation at
all. We position this paper's contribution as the causal synthesis these
four approach separately, extended along a dimension none of them test:
not only depth, position, and content-identity within one fixed system,
but content-identity specifically, carried forward through increasingly
realistic delegation structures, culminating in one the model
constructs for itself.

\section{Related Work}
\label{sec:related}

\paragraph{Cross-agent KV-cache relay.} Relaying key-value cache instead
of decoded text between agents has been proposed to avoid redundant
prefill computation. KVCOMM \citep{kvcomm2025} introduced an anchor-based
framework that estimates and corrects KV-cache ``offset drift'' when
reusing cache across different prefix contexts, reporting over 70\% cache
reuse and up to 7.8$\times$ prefill speedup on retrieval-augmented
generation, math reasoning, and coding tasks. LatentMAS
\citep{latentmas2026} extends this to full latent-space collaboration,
exchanging autoregressively-generated latent ``thoughts'' via KV-cache
working memory rather than relaying already-decoded text, reporting token
savings and accuracy gains across nine math/science/code benchmarks. We
differ from both in relaying real, decoded reasoning text through a fixed
sequential three-agent chain (Reasoner $\rightarrow$ Verifier
$\rightarrow$ Finalizer) rather than offset-corrected reuse or latent
thought generation, and in evaluating on multi-hop QA
\citep{hotpotqa2018}, which neither tested on.

\paragraph{Compressing the relayed cache.} Reducing what must be
transmitted between agents is a natural complement to relay itself. We
adopt two families of technique: layer-wise selection, which drops
calibration-identified low-importance transformer layers before
transmission and reconstructs them at the receiver via \texttt{zeros},
\texttt{nearest}, or \texttt{interpolate}, and per-head quantization,
which reduces numeric precision. A third reconstruction strategy,
Orthogonal Backfill \citep{orthobackfill2026}, injects a low-rank residual
of discarded content orthogonal to what is retained rather than
substituting or discarding it outright; we considered but did not
implement it (see Limitations for why). KV-cache compression is also
under active industrial development: KVTC \citep{kvtc2026} applies
media-compression-style transform coding -- PCA decorrelation, adaptive
quantization, entropy coding -- for up to 20$\times$ compression, a third
paradigm distinct from uniform quantization (ours, KIVI) and rotation
(TurboQuant/PolarQuant); NVIDIA's kvpress library and TensorRT-LLM's
reuse-aware cache confirm this is an actively contested axis of real
deployed systems, not a narrow academic concern.

\paragraph{Rotation-based quantization interacts badly with RoPE-encoded
keys -- a finding, not just an implementation note.} We initially
implemented a rotation-based quantization scheme in the style of
TurboQuant \citep{turboquant2026} and PolarQuant \citep{polarquant2026} --
two distinct papers sharing only two authors, not one work under two
names -- applying a fixed orthonormal Hadamard rotation to both keys and
values before quantizing to redistribute per-head outlier magnitude
(covering PolarQuant's rotate-then-quantize step, not TurboQuant's
separate QJL bias-correction term). On our real model this produced
severely out-of-distribution decoded output, not the graceful degradation
a generic rotation-based scheme predicts. We traced this to keys' rotary
position embeddings (RoPE): K is cached \emph{after} RoPE is applied,
which mixes channel pairs by a position-dependent phase, and a second,
RoPE-agnostic rotation on top disrupts that structure rather than merely
redistributing outlier magnitude. An isolating ablation (rotating only V,
leaving K unrotated) reverted the failure to ordinary quantization-noise
degradation, confirming the interaction was with K specifically --
consistent with an established interaction in the KV-quantization
literature (next paragraph), and reported here as a diagnosed negative
result, not an implementation detail.

\paragraph{RoPE-aware KV cache quantization.} A separate line of work
addresses exactly this interaction directly. KIVI \citep{kivi2024}
quantizes the key cache per-channel and the value cache per-token --
motivated by RoPE leaving channel-wise magnitude more consistent across
positions than uniform per-head ranges assume -- rather than by rotating.
RotateKV \citep{rotatekv2025} instead applies rotation \emph{before} RoPE
via Pre-RoPE Grouped-Head Rotation, explicitly designed to avoid the
interaction we observed. KVQuant \citep{kvquant2024} and related work
similarly identify post-RoPE key quantization as harder than value
quantization for exactly this reason. We adopt a scoped, KIVI-inspired fix
(per-channel key quantization only, no group-wise windowing or streaming
residual buffer from the original method) and find it resolves the
collapse -- 0.0\%/0.0\% to 52.0\%/64.1\% F1 at $n=100$, statistically
indistinguishable from our strongest layer-selection config (\texttt{D};
McNemar $p=0.86$) at higher compression (3.96x vs 3.82x, both vs.
uncompressed relay) with no calibration profile required -- direct
empirical confirmation, on our pipeline and task, of what this literature
identifies as the correct
axis for addressing the problem, rather than the rotation-based axis we
tried first. Whether KIVI's other asymmetric half (per-token value
quantization) adds anything on top is tested in Results, Finding 4.

\paragraph{Auditing whether KV reuse does what it claims.} A recent line
of work interrogates cross-agent KV/latent reuse critically rather than
only reporting end-task accuracy. ``When KV Cache Reuse Fails in
Multi-Agent Systems'' \citep{kvreusefails2026} shows reuse strategies
effective for generation agents can silently corrupt an LLM judge's
cross-candidate comparison even when end-task accuracy appears stable.
``When Does Latent Communication Pay?'' \citep{latentcommpay2026} proposes
replacing relayed KV with mismatched, zeroed, or moment-matched-random
substitutes to test whether accuracy is causally attributable to the
specific transmitted content rather than to the receiver simply having a
non-empty cache -- a methodology we adopt directly for our own causal
audit (Section~\ref{sec:method}). We extend this tradition along two axes
neither paper covers: cross-architecture generalization, treating ``does
the same technique, faithfully reproduced, transfer to a different model
family'' as itself a target of audit; and, more centrally to this paper,
delegation-structure generalization -- whether the same audit still
finds real causal content-dependence when the multi-agent topology
itself is no longer fixed in advance. Notably, \citet{latentcommpay2026}'s
own headline finding is that content-specificity is \emph{system}-dependent
-- ceiling on LatentMAS, partial on KVComm, no detected advantage on C2C,
holding audit methodology fixed and varying which published relay system
is tested. Our Finding~5's extensions to Qwen3-8B and to the C2C bridge
itself (Section~\ref{sec:results}) independently arrive at the same
qualitative pattern -- content-specificity is not universal -- while
holding audit methodology and, in the Qwen3-8B case, the relay
\emph{mechanism} fixed and varying the \emph{model} instead. Two
independent audits converging on ``it depends, and not on accuracy
alone'' from different varied axes is, we think, more informative than
either result read in isolation.

\paragraph{Model-driven delegation and recursive relay.} Recursive
Language Models \citep{rlm2026} let a root model decompose and delegate
via a live, sandboxed Python REPL rather than a topology the system
designer fixes in advance -- the model decides what to send to a sub-call
and when, and reads the sub-call's answer back as a decoded-text REPL
variable it may or may not inspect before continuing. LatentMAS
\citep{latentmas2026} and CanonicalMerge \citep{canonicalmerge2026},
discussed above, both relay something richer than text between agents,
but both over a topology fixed by the system designer -- a sequential
chain of latent thoughts and a two-child fan-in merge, respectively --
not one the model constructs for itself turn by turn. CanonicalMerge in
particular is close to our own static fan-in topology (Section
\ref{sec:method}): near-identical models (Qwen3-1.7B/4B vs.\ our
Qwen2.5-7B-Instruct/Mistral-7B-Instruct-v0.3) and a closely related task
(HotpotQA bridge questions vs.\ our full distractor-setting HotpotQA),
addressing a real methodological question we do not -- whether the
\emph{order} children are merged in biases the result -- but reporting
only downstream accuracy, with no causal-intervention test of whether
the merged content is specifically, not just numerously, necessary. We
build a genuinely dynamic, model-driven delegation loop in RLM's own
style, substitute real KV cache for its text return channel at the
sub-call boundary, and apply the same causal-audit methodology used
throughout this paper to it. To our knowledge no prior work -- RLM's own
released implementation, LatentMAS, CanonicalMerge, or the causal-audit
literature discussed above -- tests whether relayed content is causally
necessary once the delegation structure itself is no longer fixed by the
system designer.

\paragraph{Layer redundancy is not universal.} ``No Free Swap''
\citep{nofreeswap2026} finds that which transformer layers are
functionally interchangeable across Qwen3-8B, Llama-3.1-8B, and
Mistral-7B-v0.1 depends on the model and evaluation protocol rather than
being a fixed architectural property. This is consistent with our own
finding (Section~\ref{sec:results}, Finding 3) that a near-identical
proportion of dropped layers costs substantially more accuracy on
Mistral-7B-Instruct-v0.3 than on Qwen2.5-7B-Instruct, and offers one
candidate mechanism for why.

\paragraph{Three 2026 papers anticipate pieces of this paper's argument.}
``Rethinking Layer Redundancy'' \citep{rethinkredundancy2026} argues
redundancy is a joint function of model \emph{and calibration objective},
and that ``a universal layer ranking may not exist'' -- raising the
question of whether our cross-architecture gap is really about the model,
or about our calibration ranking failing to transfer. We tested the most
direct proxy: whether the calibration signal's own confidence (the
separation between kept- and dropped-layer importance scores) predicts
which model is safe to compress. It does not, and fails in the
\emph{wrong} direction -- Mistral's tier separation (0.673) is larger, not
smaller, than Qwen's (0.491), despite Mistral being the more fragile model
when that ranking is acted upon. This is a direct empirical test of
exactly the concern that paper raises, and supports a stronger reading:
not just that redundancy is calibration-objective-dependent, but that the
objective's own confidence is not a reliable signal of when its ranking is
safe to trust -- consistent with our broader claim that non-exchangeability
is not visible from a technique's internal signals.

AdaK \citep{adak2026} reports that adaptive KV-budget estimation
generalizes across Qwen3 (4B/8B) and Mistral-7B-Instruct-v0.2 -- the same
two architecture families this paper tests, not a third (Llama is not
among AdaK's evaluated models, though a superficial reading of the
abstract could suggest otherwise) -- which on its face looks like a
counterexample to our architecture-dependence claim. We do not believe it
is: AdaK estimates budget per-instance and per-model at inference time
rather than committing to a single calibration-time ranking applied
uniformly thereafter, so it never assumes a fixed, transferable ranking is
safe to act on. Read this way, AdaK's success is a positive instance of
our principle, not a counterexample -- it works specifically because it
avoids the fixed-ranking assumption our own \texttt{D} configuration
makes. This is our reading of the mechanism, not something we verified by
reimplementing AdaK ourselves; a direct head-to-head comparison on our own
pipeline is noted as a natural extension in Limitations.

Finally, ``The Pitfalls of KV Cache Compression'' \citep{pitfallskv2026}
makes a general version of the point this paper investigates concretely:
that compression methods optimize for aggregate metrics while implicitly
equating token retention with functional preservation. Our contribution is
not the abstract observation but a causally-grounded, cross-axis
demonstration of it within one system -- the same failure pattern
recurring across three structurally distinct axes of KV-cache structure
(depth, position, content-identity) in a single pipeline, with the
content-identity axis causally validated via our audit rather than
inferred from compression metrics alone. We position this paper as the
causal, cross-axis synthesis these three papers each gesture toward from a
different angle.

\section{Method}
\label{sec:method}

\paragraph{Pipeline.} We study a sequential three-agent pipeline --
Reasoner, Verifier, Finalizer -- answering HotpotQA
\citep{hotpotqa2018} (distractor configuration) questions, where each item
bundles its own ten candidate passages and requires no retrieval step;
Section~\ref{sec:results}'s GSM8K \citep{gsm8k2021} generalization check
uses the same pipeline structure. Agents share the same prompt structure
and role definitions across every configuration in this paper; the only
thing that varies between configurations is \emph{how} one agent's output
reaches the next. In \texttt{text\_agent}, the Reasoner's decoded text is
passed to the Verifier as a literal string, which reads and re-processes
it from scratch. In every KV-relay configuration (\texttt{A} through
\texttt{E}), the Reasoner's key-value cache -- its internal representation
of the question, context, and its own generated reasoning -- is passed
directly to the Verifier, which continues generation from that cache
instead of re-reading text. \texttt{single\_agent} is a no-relay control:
one model call, no intermediate agents at all. All KV-relay configurations
share identical prompts with \texttt{text\_agent}, so any accuracy or
latency delta between a KV configuration and \texttt{text\_agent} is
attributable to the relay mechanism itself, not to prompt differences.

\paragraph{Models.} We evaluate on Qwen2.5-7B-Instruct
\citep{qwen2p5_2024} (bf16) as our primary model, and -- specifically to
test whether findings generalize across architecture, not just within one
model family -- Mistral-7B-Instruct-v0.3 \citep{mistral7b2023} (bf16),
chosen for its architectural differences (distinct layer count, attention
head configuration, and tokenizer) at comparable parameter scale. Both are
decoder-only transformers using grouped-query attention and RoPE. A third
model, Qwen3-8B (Qwen Team, 2025, bf16), is evaluated on configurations
\texttt{A} and \texttt{D} only ($n=50$) specifically to extend the
architecture-dependence check in Finding 3 and the correlation in
Finding 6 to a third data point, not to replicate the full causal-audit
suite. A fourth candidate, Phi-3.5-mini-instruct, was evaluated and
dropped after isolating a pre-existing LongRoPE limitation unrelated to
our pipeline (see Limitations).

\paragraph{Compression configurations.} \texttt{B\_int8}/\texttt{B\_int4}
apply uniform per-head min-max quantization to every transmitted layer's
key and value tensors (a fix to \texttt{B\_int4}'s originally-catastrophic
failure is detailed below). \texttt{C} applies calibration-driven layer
selection -- dropping transformer layers a per-model calibration pass
identifies as low-importance before transmission, reconstructing dropped
layers at the receiver via one of three strategies (\texttt{zeros},
\texttt{nearest}, \texttt{interpolate}). \texttt{D} combines layer
selection with adaptive per-layer-tier quantization (high-importance
layers at 8-bit, medium-importance at 4-bit). \texttt{E} adds an
anchor-based cross-agent offset-correction mechanism adapted from KVCOMM
(Finding 2 is a negative result for this configuration).

\paragraph{RoPE-aware quantization fix for \texttt{B\_int4}.}
\texttt{B\_int4} (uniform 4-bit, no layer selection) initially collapsed
to 0.0\% accuracy -- diagnosed as a quantization/RoPE interaction and
confirmed via two ablations (see Related Work for the full diagnostic
story). Our fix (\texttt{B\_int4\_kivi}), inspired by KIVI, quantizes keys
per-channel (one numeric range per head dimension, computed across
positions) instead of per-head, leaving values unchanged.

\paragraph{Statistical methodology.} For every headline comparison between
two configurations, we report exact (binomial-form) McNemar's test on
paired per-example correctness -- appropriate because both configurations
are evaluated on the identical example set -- and a paired bootstrap 95\%
confidence interval (10,000 resamples) for the difference in mean
token-F1. We additionally report Monte Carlo-estimated statistical power
at the current sample size for comparisons that do not reach significance,
so that ``not significant'' and ``confirmed no effect'' are never
conflated in our reporting.

\paragraph{Causal audit.} To test whether a configuration's accuracy
reflects the receiving agent using the specific transmitted content,
rather than merely having \emph{some} non-empty cache to attend over, we
substitute the relayed KV cache with (i) an all-zero cache, (ii) a
moment-matched random cache (same per-tensor mean/standard deviation, no
real structure), and (iii) another, unrelated held-out question's real
cache at the same pipeline stage (sampled from a pool built from examples
never included in the scored evaluation set, so no scored question's
content can leak back to itself). We apply this to our three strongest
configurations (\texttt{A}, \texttt{B\_int8}, \texttt{D}).

We deliberately do not treat these three conditions as interchangeable
evidence for the same claim. Zero-ablation is a known-imperfect causal
baseline in the interpretability literature: it pushes activations off the
training distribution, which can produce large behavioral changes for
reasons unrelated to whether the zeroed content carried meaningful
information. We treat the moment-matched-random and mismatched-example
conditions -- which preserve realistic activation statistics or substitute
another real, in-distribution cache -- as the more informative evidence,
and zero-ablation as a supplementary sanity floor; if zeroing alone
diverges sharply from the other two, we attribute that to the
out-of-distribution effect rather than content-dependence. A fully
rigorous treatment of ``no difference'' would also use a pre-registered
equivalence margin rather than absence of significance under our paired
test; we did not have the budget to pre-register and power such a test,
and report significance/CI results with that caveat.

\paragraph{Fan-in topology.} To test whether content-identity
non-exchangeability is an artifact of our sequential chain specifically,
we build a static fan-in decomposition: HotpotQA's ten candidate passages
are split across two independent child agents, each computing its own KV
cache over its half; the two caches are RoPE-position-shifted (reusing
the same primitive as our offset-correction baseline) and concatenated
before one aggregator attends over both and answers, so the aggregator
receives two parallel sources rather than one sequential handoff. The
causal audit substitutes \emph{every} child's KV (not a single child,
holding the other's real content available as an unaudited shortcut --
an early design bug we found and fixed, see Limitations), matching the
full-relay-corruption strength of the sequential pipeline's own audit.

\paragraph{Dynamic, model-driven delegation (RLM+KV).} Our central
topology extension builds a real, sandboxed Python REPL loop in RLM's
own style \citep{rlm2026}: the root model writes and executes code
turn by turn, deciding for itself which passages to inspect directly and
when to delegate a sub-question to a child call, continuing until it
calls a terminal \texttt{final\_answer(...)}. We implement two
sub-call return channels to isolate the variable RLM itself does not
test: \texttt{text}, matching RLM's own mechanism, decodes a sub-call's
answer to text and feeds it back as new tokens; \texttt{kv} instead
RoPE-shifts and splices the sub-call's own real KV cache directly onto
the root's running cache, so the root never observes the sub-call's
answer as text at all. Both conditions require the root to carry a
real, continuously-extended KV cache across turns rather than
re-tokenizing the full transcript each turn, unlike RLM's own reference
implementation, so that the two return channels differ only in what is
relayed, not in how the root's own generation mechanism works. The
causal audit substitutes a sub-call's real KV with zeroed, random, or a
held-out question's real content immediately before the splice,
identically to the two fixed topologies. One structural difference from
both fixed topologies is a genuine limitation, not an oversight: the
root retains direct sandboxed access to the original passages
independent of any sub-call, so corrupting a sub-call's relayed KV does
not remove all of the root's information the way corrupting an entire
relayed context does in the fixed topologies (see Results and
Limitations for how this shapes what the audit can and cannot show).

We implement this ourselves directly on \texttt{transformers}
\citep{transformers2020} rather than on top of \texttt{dspy.RLM}, the
first-party DSPy implementation of RLM by one of the paper's own
authors. This is a structural necessity, not a preference: DSPy is
built on \texttt{litellm} and is, by design, a text-only, API-centric
orchestration layer -- a sub-call through \texttt{dspy.RLM}'s own
\texttt{llm\_query} returns decoded text and nothing else, with no
exposure anywhere in the framework to a model's KV cache or other
internal state. Our \texttt{kv} return channel is not implementable on
top of it. We verified this directly against DSPy's own documentation
rather than assume it from the framework's general reputation before
committing to a from-scratch implementation.

\section{Results}
\label{sec:results}

Eight findings follow. Two are the paper's headline results: Finding 5,
the causal audit establishing that relayed KV carries specific, real
content rather than merely being non-empty, and Finding 6, that a
calibration procedure's own confidence fails to predict -- and actively
misdirects on -- which architecture is safe to compress. The remaining
findings characterize specific published techniques against this
backdrop.

\begin{table*}[t]
\centering
\small
\begin{tabular}{lcccc}
\toprule
\textbf{Config} & \textbf{Accuracy} & \textbf{F1} & \textbf{Latency} & \textbf{KV/hop} \\
\midrule
single\_agent & 57.0\% & 68.4\% & 1.6s & --- \\
text\_agent & 54.0\% & 68.3\% & 7.2s & --- \\
A & 56.0\% & 69.1\% & 8.6s & 144.35 MB \\
B\_int8 & 57.0\% & 70.5\% & 8.3s & 72.16 MB (2.00x) \\
B\_int4 (original) & 0.0\% & 0.0\% & 39.4s & 38.38 MB (4.00x) \\
\textbf{B\_int4\_kivi (fixed)} & \textbf{52.0\%} & \textbf{64.1\%} & 9.6s & 36.43 MB (3.96x vs A) \\
C & 43.0\% & 57.8\% & 10.0s & 104.16 MB \\
D & 50.0\% & 61.9\% & 10.2s & 37.82 MB (3.82x vs A) \\
E & 9.0\% & 17.4\% & 19.8s & 35.63 MB \\
E\_int8 & 1.0\% & 6.2\% & 22.2s & 49.05 MB \\
\bottomrule
\end{tabular}
\caption{Established results (Qwen2.5-7B-Instruct, HotpotQA, $n=100$).
Ratios are compressed-vs-A's 144.35~MB, except B\_int4 and B\_int8, whose
self-referential ratio equals their vs.-A ratio exactly since neither
drops any layers. See Method for the real nibble-packing that makes
4-bit configs genuinely, not just nominally, smaller than 8-bit ones.}
\label{tab:main}
\end{table*}

\begin{figure}[t]
\centering
\includegraphics[width=\columnwidth]{figures/fig3_compression_pareto.png}
\caption{Compression vs. accuracy across every relay condition in
Table~\ref{tab:main}. \texttt{B\_int4}'s collapse shows smaller is not
automatically better; \texttt{B\_int4\_kivi} is the Pareto-best point,
matching \texttt{D}'s accuracy at the smallest relayed cache of any
configuration tested.}
\label{fig:compression_pareto}
\end{figure}

\paragraph{Finding 1 -- uniform 8-bit KV quantization is statistically
indistinguishable from uncompressed relay.} \texttt{A} vs \texttt{B\_int8}:
McNemar $p=1.0$, only 3 discordant examples out of 100. Confirmed on a
second model family (Mistral-7B-Instruct-v0.3, $n=50$): 0 discordant
examples -- identical per-example outcomes. This is not an underpowered
null result; the observed effect is tight enough at this sample size to be
a genuine near-zero cost.

\paragraph{Finding 2 -- KVCOMM-style offset correction (\texttt{E}) fails
catastrophically, independent of a decoding confound we explicitly ruled
out.} \texttt{D} vs \texttt{E}: McNemar $p\approx0.00003$ (both before and
after an unrelated repetition-penalty fix that changed several other
configurations' numbers -- see Limitations). Manual inspection of
\texttt{E}'s raw output shows duplicated phrases and hallucinated content
inconsistent with the anchor-corrected cache carrying usable information,
not merely a lower-quality-but-coherent answer.

\paragraph{Finding 3 -- layer-selection's accuracy cost is
architecture-dependent.} On Mistral-7B-Instruct-v0.3 ($n=50$), dropping
9/32 layers (28\%) cost 30.5 F1 points relative to uncompressed relay
($p=0.0001$). On Qwen2.5-7B-Instruct ($n=100$), a near-identical
proportion (8/28, 29\%) shows a smaller, not-yet-significant cost
($p=0.34$) -- the cross-model claim does not rest on that reaching
significance, though: even Qwen's most generous uncertainty bound (F1 CI
upper +16.4 points) is under half of Mistral's confirmed collapse. A
third model, Qwen3-8B ($n=50$), sharpens rather than complicates this
contrast: \texttt{D} numerically edges out \texttt{A} (56.0\% vs.\
54.0\%), but this is noise, not a real effect -- McNemar $p=1.0$ on only 3
discordant pairs out of 50, F1 bootstrap CI crossing zero. Read
correctly, this is a second model, after Qwen2.5, where layer-selection's
cost is statistically indistinguishable from zero, against Mistral's
single confirmed, significant collapse -- evidence for
architecture-dependence specifically, not for ``layer-selection is
usually costly'' or ``usually free.'' Consistent with prior work finding
layer redundancy architecture- and protocol-dependent, not a fixed
property of a given compression ratio.

\paragraph{Finding 4 -- \texttt{B\_int4}'s fixed variant matches our
strongest layer-selection configuration's accuracy at genuinely higher
compression and lower complexity.} \texttt{B\_int4\_kivi}
(52.0\%/64.1\%) vs \texttt{D} (50.0\%/61.9\%): McNemar $p=0.86$, F1 delta
+2.2 points [95\% CI: $-8.2$, $+12.8$] -- statistically indistinguishable
-- while achieving 3.96x compression versus \texttt{D}'s 3.82x (both vs.\
uncompressed \texttt{A}), and requiring no per-model calibration profile.
This number went through a real correction cycle we report for
transparency: this codebase's quantized tensors were initially stored as
one full \texttt{torch.uint8} per element regardless of bit-width -- no
nibble-packing existed anywhere in the compressor, so a 4-bit value took
the same byte an 8-bit one did, and our first-reported ratios (silently
assuming packing that did not exist) were wrong in a way that happened to
be numerically optimistic. We caught this, then rather than only
correcting the byte \emph{accounting}, implemented real
two-values-per-byte nibble packing (unit-tested: round-trip pack/unpack
exactness, and \texttt{compress}/\texttt{decompress} bit-identical output
with and without packing -- packing is a storage-format change, not a
numeric one) so the reported ratio reflects genuine physical storage. The
corrected, packing-backed numbers land back at the originally-reported
figures almost exactly, because the original formula's \emph{arithmetic}
(0.5 bytes/element for 4-bit) was always what real packing would produce
-- the bug was that nothing in the compressor actually packed anything.
\texttt{B\_int4\_kivi} also trends below \texttt{A}/\texttt{B\_int8} on
accuracy (52\% vs 56--57\%), but this gap is not yet statistically
confirmed at $n=100$ (7--12 discordant examples, $p=0.36$--$0.50$); we
report this as an open question, not as cost-free the way Finding~1
establishes for \texttt{B\_int8}. Also fixing the value side
(\texttt{B\_int4\_kivi\_full}, per-token quantization) does not improve
on this further (44.0\% vs 50.0\% on a matched 50-example subset,
$p=0.51$, trending the other way) and is still worse on compression
(3.75x vs \texttt{B\_int4\_kivi}'s 3.99x self-referential) for a real,
mechanical reason: per-token quantization stores a scale/zero-point pair
per sequence position, and sequence length (hundreds to low thousands of
tokens) vastly exceeds K's per-channel grouping (128 channels) -- real
overhead, no accuracy benefit, and packing does not remove it. Together
with an earlier null result for rotating values instead, two
independently-motivated value-side interventions both failed -- evidence
the key/RoPE interaction was the entire mechanism behind
\texttt{B\_int4}'s collapse, not one of several factors.

\begin{figure}[t]
\centering
\includegraphics[width=\columnwidth]{figures/fig1_causal_audit.png}
\caption{Causal audit: real $>$ mismatched $>$ zeroed $\approx$ random,
replicated across \texttt{A}, \texttt{D}, and \texttt{B\_int8}. Every
pairwise comparison is statistically significant.}
\label{fig:causal_audit}
\end{figure}

\paragraph{Finding 5 -- headline result: relayed KV demonstrably carries
specific, real content, not merely a non-empty cache.} We ran the causal
audit (Section~\ref{sec:method}) on \texttt{A} (uncompressed relay,
sequential chain), $n=50$: 50.0\%/64.0\% real vs.\ zeroed/random both
0.0\%/0.0\% vs.\ mismatched 28.0\%/37.1\%, every pairwise comparison
significant (real vs.\ zeroed/random $p<0.0001$ each; real vs.\
mismatched $p=0.0127$; mismatched vs.\ zeroed/random $p=0.0001$ each).
This is, to our knowledge, the first direct causal validation of this
kind for a sequential multi-agent KV-relay chain, and it underwrites
every compression result in this paper: Findings 1 and 4 are only
meaningful claims about \emph{preserved information} if the underlying
relay is shown to carry real content in the first place, which this
finding establishes directly rather than assuming. Qualitatively,
\texttt{mismatched}'s wrong answers are coherent and grammatical -- real
linguistic structure attached to the wrong question -- \texttt{zeroed}
produces garbled but recognizable English, and \texttt{random} is
\emph{more} degraded still (code-fragment tokens, mixed-language noise)
despite matching real activation statistics -- we read this as a zero
vector being a weak key attention can down-weight, while realistic
random noise resembles genuine signal and actively misdirects it.

Table~\ref{tab:causal_audit_full} reports every subsequent extension of
this audit in one place, rather than as a chain of similar-sounding
paragraphs: five relay configurations on Qwen2.5-7B-Instruct/HotpotQA
(the remaining rows in the top block), five on
Mistral-7B-Instruct-v0.3, two on Qwen3-8B, one on GSM8K, and, central to
this paper, one each on the two topology extensions discussed below.
Across every Qwen2.5/HotpotQA configuration (\texttt{D}, \texttt{B\_int8},
\texttt{C}, \texttt{B\_int4\_kivi}), the same three-tier ordering
replicates with real vs.\ zeroed/random significant at $p<0.0001$ in
every case and real vs.\ mismatched significant throughout
($p=0.0013$--$0.0192$); a single GSM8K \citep{gsm8k2021} confirmatory
run replicates it more decisively still ($p<0.0002$ throughout), showing
the claim is not an artifact of one task. \texttt{Mismatched} vs.\
\texttt{zeroed}/\texttt{random} across all these configurations lands at
$p=0.0001$--$0.0005$, tighter than \texttt{A}'s own $p=0.0001$.

% Consolidated causal-audit table -- draft for review before insertion into main.tex.
% Covers every real/zeroed/random/mismatched comparison currently spread
% across six "Finding 5, extended..." paragraphs, plus the new topology
% generalization rows (fan-in, RLM+KV) not yet in main.tex.

\begin{table*}[t]
\centering
\small
\begin{tabular}{llcccccc}
\toprule
\textbf{Config} & \textbf{Model / Task / Topology} & \textbf{Real} & \textbf{Zeroed} & \textbf{Random} & \textbf{Mismatched} & \textbf{$n$} \\
\midrule
\multicolumn{7}{l}{\emph{Qwen2.5-7B-Instruct, HotpotQA, sequential chain}} \\
A            & --                          & 50.0/64.0 & 0.0/0.0  & 0.0/0.0  & 28.0/37.1 & 50 \\
D            & --                          & 54.0/60.2 & 0.0/0.0  & 0.0/0.2  & 26.0/32.2 & 50 \\
B\_int8      & --                          & 54.0/68.0 & 0.0/0.0  & 0.0/0.0  & 24.0/34.4 & 50 \\
C            & --                          & 46.0/55.7 & 0.0/0.0  & 0.0/0.0  & 24.0/38.9 & 50 \\
B\_int4\_kivi& --                          & 50.0/61.7 & 0.0/0.0  & 0.0/0.0  & 28.0/34.7 & 50 \\
\midrule
\multicolumn{7}{l}{\emph{Generalization: model, task, and topology}} \\
A            & Mistral-7B-Instruct-v0.3    & 56.0/63.3 & 0.0/0.5  & 0.0/0.3  & 16.0/26.4 & 50 \\
C            & Mistral-7B-Instruct-v0.3    & 20.0/30.3 & 0.0/0.5  & 0.0/0.1  & 14.0/26.1 & 50 \\
D            & Mistral-7B-Instruct-v0.3    & 22.0/32.8 & 0.0/0.5  & 0.0/0.1  & 16.0/29.3 & 50 \\
B\_int8      & Mistral-7B-Instruct-v0.3    & 56.0/63.3 & 0.0/0.5  & 0.0/0.1  & 20.0/29.3 & 50 \\
B\_int4\_kivi& Mistral-7B-Instruct-v0.3    & 54.0/63.2 & 0.0/0.5  & 0.0/0.1  & 16.0/25.0 & 50 \\
A            & Qwen3-8B                    & 54.0/67.2 & 46.0/55.8& 0.0/0.0  & 64.0/74.8 & 50 \\
D            & Qwen3-8B                    & 56.0/71.2 & 46.0/55.8& 2.0/5.0  & 52.0/62.4 & 50 \\
A            & GSM8K (Qwen2.5)             & 90.0/---  & 0.0/---  & 2.0/---  & 28.0/---  & 50 \\
kv (fan-in)  & Fan-in topology (Qwen2.5)   & 40.0/53.0 & 0.0/0.8  & 0.0/0.0  & 6.0/12.0  & 50 \\
kv (RLM+KV)  & Dynamic delegation (Qwen2.5)& 24.0/32.3 & 8.0/11.2 & 6.0/9.2  & 10.0/15.5 & 50 \\
\bottomrule
\end{tabular}
\caption{Every real/zeroed/random/mismatched causal-audit comparison in this
paper: accuracy/F1 (\%) per condition. \emph{Real} vs.\ \emph{zeroed}/\emph{random}
is significant ($p<0.0001$) in every row except Qwen3-8B, where content-dependence
itself is architecture-dependent (see text). \emph{Real} vs.\ \emph{mismatched}
is significant in all rows except Mistral \texttt{C}/\texttt{D} (underpowered on
an already-degraded layer-selection baseline, see Finding 3) and Qwen3-8B (a
distinct, architecture-dependent finding, see text). GSM8K reports accuracy only
(no partial-credit F1 metric for this task).}
\label{tab:causal_audit_full}
\end{table*}


On Mistral, \texttt{B\_int4\_kivi} replicates the full five-comparison
ladder as cleanly as on Qwen, but \texttt{C} tells a more nuanced story:
mismatched clearly beats zeroed/random ($p=0.0156$ each, content still
matters at all), but \texttt{C} itself is not statistically
distinguishable from mismatched ($p=0.51$, F1 CI includes zero) --
unlike on Qwen, where this comparison was significant. The most likely
explanation is not that content-identity stops mattering on Mistral, but
that \texttt{C} already performs poorly there before any auditing
(20.0\%, well below Qwen's 46.0\%), an independent confirmation of
Finding 3's claim that Mistral is unusually fragile to layer-selection
specifically -- less headroom to detect a further gap on an
already-degraded baseline, not a failure to replicate. \texttt{D} on
Mistral shows the identical partial pattern (mismatched vs.\
zeroed/random significant, $p=0.0078$ each; real vs.\ mismatched not,
$p=0.375$), a second, independent confirmation of the same explanation.
\texttt{B\_int8} on Mistral replicates the full clean pattern, matching
every quantization-only configuration tested; the dividing line across
two separate experiments is now consistent: layer-selection configurations
(\texttt{C}, \texttt{D}) show the attenuated pattern on Mistral,
quantization-only configurations (\texttt{A}, \texttt{B\_int8},
\texttt{B\_int4\_kivi}) show the full pattern throughout.

On Qwen3-8B, content-dependence itself turns out to be
architecture-dependent, a genuinely different pattern from either prior
model, not a weaker version of the same one: random substitution still
causes total, highly significant collapse ($p<0.0001$ on both
\texttt{A} and \texttt{D}), but real is \textbf{not} statistically
distinguishable from either zeroed or mismatched (all four comparisons
$p=0.30$--$0.79$, F1 CIs include zero). We verified this is a real
effect, not a bug, three ways: no donor-question pool leakage;
\texttt{A\_audit\_zeroed} and \texttt{D\_audit\_zeroed} produce
byte-identical output on all 50 examples, which looked alarming until we
traced it to \texttt{D}'s own \texttt{zeros} reconstruction strategy
composing correctly with the audit's own zeroing, not a bug; and real,
zeroed, and mismatched overlap on only 10--24 of 50 examples pairwise,
confirming the substitution takes effect selectively. Read together with
Finding 6, this extends the same meta-point to the content-identity
axis: even this paper's own central mechanism is not a universal
property of KV relay, but itself varies by architecture, with a
stronger/larger model here robust enough to partially recover from
losing or misdirecting real content while remaining just as defenseless
against structural noise.

Everything above is on our fixed sequential chain. Figure~\ref{fig:topology_ladder}
previews where the rest of Results is headed: the same three-tier
ordering, panel by panel, across the two topology extensions below --
this paper's central contribution.

\begin{figure}[t]
\centering
\includegraphics[width=\columnwidth]{figures/fig4_topology_ladder.png}
\caption{The causal-audit ladder across all three topologies tested in
this paper. Real content beats mismatched content, which beats
zeroed/random, in every topology -- including one the model constructs
for itself turn by turn, not only a fixed pipeline. The recursive
topology's non-zero zeroed/random floor is a real, discussed scope
difference (Section~\ref{sec:method}), not noise.}
\label{fig:topology_ladder}
\end{figure}

\paragraph{Finding 5, extended to a second topology -- the cleanest
three-tier result in this paper, and content-dependence is not specific
to a sequential chain.} We built a static fan-in decomposition
(Section~\ref{sec:method}): HotpotQA's ten passages split across two
independent child agents, merged via RoPE-shifted concatenation before
one aggregator attends over both. Substituting \emph{every} child's KV
($n=50$, HotpotQA; see Method and Limitations for a real
experimental-design bug this required fixing -- substituting only one
child understated the effect by leaving the other child's real content
as an unaudited shortcut) gives the full three-tier ordering, and more
decisively than the sequential chain: 40.0\%/53.0\% real vs.\ 0.0\%/0.8\%
zeroed vs.\ 0.0\%/0.0\% random vs.\ 6.0\%/12.0\% mismatched. Real vs.\
zeroed: $p<0.0001$ (F1 delta $+0.522$ [$+0.396$, $+0.645$]). Real vs.\
random: $p<0.0001$ (F1 delta $+0.530$ [$+0.403$, $+0.654$]). Real vs.\
mismatched: $p=0.0001$, 18 of 19 discordant pairs favoring real (F1
delta $+0.410$ [$+0.267$, $+0.546$]) -- the strongest, most decisive
version of this specific comparison anywhere in this paper. Unlike the
recursive topology below, this static decomposition has no alternate
information pathway for a corrupted child's answer to route around, so
zeroed and random collapse to genuinely near-total zero, matching every
audit on the sequential chain. One precise nuance: mismatched vs.\
zeroed/random is not significant by the exact McNemar test ($p=0.25$
each, only 3 discordant pairs -- the test's own floor at that count, not
necessarily a true null), but the F1 bootstrap CI excludes zero for both
($+0.112$ [$+0.035,+0.200$] and $+0.120$ [$+0.047,+0.206$]) --
accuracy-level evidence is underpowered at this $n$ for this specific
leg, while the finer-grained F1 comparison trends real; we report both
readings rather than rounding to whichever is more convenient. A
garbage-signature scan of all 200 generations (repetition loops,
out-of-vocabulary tokens) gives a clean, monotonic qualitative match to
the accuracy numbers: 0/50 flagged for real, 4/50 mismatched, 23/50
zeroed, 48/50 random -- the same ``random is worse than zeroed'' texture
as Finding 5's original observation, replicated cleanly a second time.

\paragraph{Finding 5, extended to a third topology -- content-dependence
survives a deliberately unfavorable stress test.} We built the RLM-style
dynamic delegation loop described in Section~\ref{sec:method}: the model
itself decides, turn by turn, what to delegate and when, with a sub-call's
answer relayed as real KV cache (\texttt{kv}) rather than RLM's own
decoded-text channel (\texttt{text}). We are explicit that this is the
least favorable of our three topologies for demonstrating anything: our
distractor-setting HotpotQA questions (ten passages, a few thousand
tokens) do not need RLM-style decomposition at all -- unlike the
context regimes RLM itself targets -- so a real-model REPL loop
constructing its own delegation strategy here is solving a harder
coordination problem than the task requires, and getting it to produce
coherent output at all needed real scaffolding (Section~\ref{sec:method},
grounding and low-exploration nudges). Consistent with that framing,
$\texttt{kv}$ vs.\ $\texttt{text}$ accuracy remains statistically
unresolved at $n=50$ (24.0\%/32.3\% vs.\ 22.0\%/28.5\%, $p=1.0$) -- we do
not claim an accuracy win here, and this topology's contribution does not
rest on one: the question this section answers is not ``does KV relay
beat text relay under model-driven delegation'' but ``does content-identity
non-exchangeability survive even here, where the topology is at its most
adversarial to any clean result.'' The causal audit does resolve: $\texttt{kv}$ vs.\
$\texttt{kv\_audit\_zeroed}$, $p=0.0078$ (F1 delta $+0.210$
[$+0.110,+0.323$]); vs.\ $\texttt{kv\_audit\_random}$, $p=0.0039$ (F1
delta $+0.230$ [$+0.127,+0.347$]); and, completing the ladder, vs.\
$\texttt{kv\_audit\_mismatched}$, $p=0.0156$ (F1 delta $+0.167$
[$+0.071,+0.274$]) -- the full three-tier signature, now confirmed on a
topology the model constructs for itself rather than one fixed in
advance. One honest, mechanistically important caveat, unique to this
topology: zeroed/random do not collapse to near-zero the way every other
audit in this paper does (8.0\%/6.0\% rather than $\approx$0\%), because
here uniquely the root retains direct sandboxed access to the source
passages independent of any sub-call (Section~\ref{sec:method}) --
corrupting a sub-call's answer does not remove all of the root's
information the way corrupting an entire relayed context does elsewhere
in this paper. This audit is real and significant, but structurally
narrower than the fixed topologies' own, and we report it as such.

A second, complementary signal, mined from the same transcripts at no
additional cost, gives a more complete confirmation than accuracy alone:
the model's own \emph{behavior} -- whether it commits to an answer within
its turn budget, not only what it eventually answers -- tracks the same
ordering, including the one comparison accuracy could not resolve.
Timeout rate, mean turn count, and mean delegation-call count are all
monotonic across conditions (real 22\% timeout / 7.72 turns / 1.38 calls;
mismatched 42\%/8.48/1.12; zeroed 72\%/9.46/0.76; random 76\%/9.60/0.66).
Substituting ``session finished without exhausting its turn budget'' for
``correct'' in the same paired McNemar test, every comparison is
significant, including mismatched vs.\ random ($p<0.0001$) -- exactly
the leg final-answer accuracy left unresolved (real vs.\ mismatched
$p=0.0309$; real vs.\ zeroed/random $p<0.0001$ each; mismatched vs.\
zeroed $p=0.0001$). Even though the model never observes the corrupted
content as text, its downstream process is gradedly sensitive to content
quality in a way the coarser final-answer metric partially masks --
a genuinely new angle relative to every other causal audit in this
paper, all of which score only the final answer, and a plausible
mechanistic account of why a real content-dependence effect can survive
in behavior even where an accuracy-only view would call it unresolved.
We verified this run shows no new degenerate output at this scale (0/50
flagged for real/text; a clean, monotonic 3/50--48/50 ordering across the
remaining conditions matching the fan-in topology's own signature).

An earlier version of this run was not run-to-run deterministic: diffing
against an independent rerun on the identical 50 questions originally
found 30\% of \texttt{text}-channel raw generations differing under
greedy decoding, plausibly GPU floating-point non-associativity
cascading through this prototype's long, self-referential, code-writing
decode loop -- unlike the fixed topologies, which are bit-identical
across reruns. We forced deterministic CUDA/cuBLAS/cuDNN kernels
(\texttt{torch.use\_deterministic\_algorithms}, a fixed
\texttt{CUBLAS\_WORKSPACE\_CONFIG} set before any CUDA context exists)
and reran the full five-condition ladder independently: every reported
number -- exact match, F1, timeout rate, and delegation-call count --
reproduces bit-identically against the run above (0 of 250 raw answers
differ across all five conditions). The one residual nuance, found by
diffing every field rather than only the reported ones: 33 of the
\texttt{random} condition's 50 examples show different intermediate
turn-by-turn wording before both runs independently exhaust their turn
budget on the identical question with the identical outcome -- the
highest-entropy condition's exploratory path is not perfectly pinned
down, but no scored quantity in this paper is affected. The results
reported in this section are therefore confirmed reproducible, not a
single sample from an uncontrolled process.

\paragraph{Finding 6 -- headline result: a calibration signal's own
confidence does not predict downstream layer-selection safety, and fails
in the wrong direction.} We tested the most direct available proxy for
whether our cross-architecture gap (Finding 3) reflects the model or
merely our calibration procedure: the separation between
calibration-importance scores assigned to kept (tier-1/2) versus dropped
(tier-3) layers, as a measure of how confidently the calibration ranks
layers. Qwen's mean separation is 0.491 (tier-1 0.737, tier-3 0.246);
Mistral's is \emph{larger}, 0.673 (tier-1 0.851, tier-3 0.178) --
Mistral's calibration looks more decisive, not less, yet Mistral is the
model whose accuracy collapses harder when that ranking is acted upon
(Finding 3). This rules out a noisier or less-confident calibration signal
for Mistral as the explanation for Finding 3, and supports a sharper one:
the calibration objective can be equally or more internally confident
while being less reliable, undetectable from the calibration output
alone. Adding Qwen3-8B as a third point (separation 0.591, between
Qwen2.5's and Mistral's, pairing with \texttt{D}'s near-zero cost from
Finding 3) lets us correlate tier-separation against degradation directly
rather than compare only two models qualitatively: Pearson $r=0.619$,
Spearman $r=0.500$ across the three. We report this only as a trend
consistent with the two-model finding, not as validating evidence --
$n=3$ is barely more informative than $n=2$, and we would not cite this
correlation as meaningful without at least two or three further model
families (see Limitations for why a fourth did not make it into this
draft). This remains a real, falsified hypothesis test, not an assumed
conclusion, but not yet a validated general predictor either.

\paragraph{Finding 7 -- architecture-dependent layer-selection failures
differ in kind, not only in rate.} Beyond the accuracy gap already
reported (Finding 3), \texttt{D}'s incorrect answers differ in character
between models. Qwen's are short (mean 18 characters, 2\% exceed 150) and
overwhelmingly close, traceable near-misses. Mistral's are markedly longer
(mean 109 characters, 15\% exceed 150, maximum 1,454) and include a
genuine long tail of severe failures -- fluent but entirely fabricated
tangents, degenerate repetition loops -- absent from Qwen's failure set.
We checked an obvious alternative explanation -- under-penalizing
repetition for Mistral relative to its own tuned defaults -- against each
model's published generation defaults, and find the opposite: Qwen's own
default \emph{is} our \texttt{repetition\_penalty=1.05}, while Mistral's
specifies none, so our setting applies \emph{more} correction for
Mistral, not less. The long-tail pattern persists anyway, weakening a
decoding confound as the explanation and supporting depth-axis
non-exchangeability differing in \emph{character}, not only magnitude,
across architectures.

\begin{figure}[t]
\centering
\includegraphics[width=\columnwidth]{figures/fig2_b_int4_journey.png}
\caption{The \texttt{B\_int4} diagnosis-to-fix journey: collapse, a broken
rotation-based attempt, an isolating diagnostic, the working fix
(\texttt{B\_int4\_kivi}), and two value-side variants that did not
improve on it.}
\label{fig:b_int4_journey}
\end{figure}

\paragraph{Finding 8 -- donor-question content occasionally, but rarely,
bleeds through under mismatched substitution.} Unlike ``When Latent
Agents Lie'' \citep{latentagentslie2026}, which studies an adversarial
agent deliberately substituting hidden state in a fan-in topology, our
substitution is a non-adversarial intervention in a sequential chain: when
the receiving agent is wrong, does its answer reflect the donor's
content, or ordinary confusion on the real question? Manually reviewing
20 of 36 incorrect \texttt{A\_audit\_mismatched} answers against their
donor question, we find one unambiguous instance -- a Kansas university
fight-song question produced ``Ellie Goulding'' in its wrong answer,
traceable only to a donor question about that singer -- and one weaker,
ambiguous case; the remainder show ordinary failure on the real question's
own topic. A genuine but rare phenomenon, not the dominant explanation for
\texttt{A\_audit\_mismatched}'s accuracy loss -- a manual review of a
subset, not an exhaustive classification (see Limitations).

\section*{Limitations}

\paragraph{Causal audit scope, bleed-through, and remaining topology
questions.} The causal audit now covers every relay condition in this
paper except \texttt{B\_int4} itself, which is uninformative to audit
given it already collapses unaudited (Section~\ref{sec:results}). The
bleed-through analysis (Finding 8) remains narrower: a manual review of
20 of 36 incorrect \texttt{A\_audit\_mismatched} examples by one
annotator, with no automated classifier or inter-annotator agreement
check -- a qualitative, exploratory observation, not a validated rate; a
rigorous version would need an automated or multiply-annotated
classification over the full set of incorrect examples, which we did not
have the budget to complete. On the recursive delegation topology
specifically, two concrete engineering findings are worth stating plainly
rather than only implying from Results: getting a working, coherent
baseline required diagnosing and fixing several real failure modes (a
sub-call's own prompt framing leaking into the receiver's attention when
spliced without isolating the generated answer span; a session-wide
repetition-tracking loss specific to a persistent, turn-continuing cache;
a RoPE position-shift bug specific to splicing an answer-only, prompt-
truncated child KV, whose fix we verified two ways -- a corrected unit
test and a real-model rerun showing every previously degenerate example
become coherent); and even with those fixed, the $\texttt{kv}$-vs-$\texttt{text}$
accuracy comparison remains genuinely unresolved at $n=50$, not merely
underpowered in a way we expect more data to resolve trivially --  we
report the causal-audit result, which does resolve, as this topology's
real contribution, and leave the accuracy comparison open rather than
claim more than the evidence supports. Beyond the three topologies
tested here, a genuinely adversarial fan-in topology \citep{latentagentslie2026}
and LatentMAS's layer-wise KV concatenation \citep{latentmas2026} remain
untested against our specific causal-audit methodology; both are real,
cited, not-yet-attempted extensions, not gaps we found and
characterized without naming.

\paragraph{Single-process, single-GPU evaluation.} Every experiment runs
within one Python process on one RTX 4090: no \texttt{KVMessage} is ever
serialized or transmitted over a network. Compression-ratio and
transmitted-byte accounting reflects what a real distributed deployment
would need to send and is meaningful on that basis, but our latency
numbers exclude network transfer time, and smaller relayed caches are not
observed to be faster in this setup, since the model still executes every
transformer layer regardless of what its injected cache contains -- claims
here characterize communication-cost tradeoffs and accuracy effects, not
measured end-to-end deployment speedups. This single-GPU constraint also
bounded the sample sizes and the number of model/config combinations we
were able to test; a systems-scale deployment study (multi-GPU, real
network transfer, throughput under concurrent load) was out of scope.

\paragraph{Statistical power and replication.} We report paired
significance tests (exact McNemar) and bootstrap confidence intervals for
every headline comparison; several -- INT8's near-zero cost, offset
correction's catastrophic failure, Mistral-7B's layer-selection collapse
-- are significant at conventional thresholds even at our evaluation
scale. Others are not: at $n=100$, we cannot yet statistically distinguish
Qwen2.5-7B's layer-selection condition from its uncompressed baseline in
isolation (power $\approx$15\%; roughly 3$\times$ the current sample size
would be needed). Our cross-architecture claim rests on the independently
well-powered Mistral-7B result, not the underpowered Qwen-only
comparison. Relatedly, \texttt{B\_int4\_kivi}'s comparison against
\texttt{D} (statistically indistinguishable, McNemar $p=0.86$) rests on a
single $n=100$ run each, not a multi-seed or replicated evaluation -- we
report the confidence interval alongside the point estimate specifically
so ``matches \texttt{D}'' is read as the current best estimate under one
run, not a claim strengthened by replication we have not performed.

\paragraph{Scope: three models (one only partially), mostly one task,
greedy decoding.} We test Qwen2.5-7B-Instruct and Mistral-7B-Instruct-v0.3
(chosen for architectural diversity at comparable scale) on the full
suite of configurations, and Qwen3-8B on \texttt{A}/\texttt{D} only
($n=50$), specifically to extend Finding 3's architecture-dependence
claim and Finding 6's tier-separation correlation to a third point -- not
to replicate the causal audit or compression-fix findings on a third
architecture. A fourth candidate, Phi-3.5-mini-instruct, was evaluated
and dropped: after fixing two real bugs in how our pipeline hands decode
off to \texttt{model.generate()} under non-default RoPE scaling, the
model still failed to produce coherent output once the pipeline's
cumulative relayed context crossed roughly 4{,}096 tokens -- the model's
\texttt{original\_max\_position\_embeddings}. We isolated this to a
pre-existing limitation of the model/library combination itself, not our
pipeline: plain, unmodified \texttt{model.generate()}, with no relay,
injection, or pipeline code involved at all, degrades into the identical
repetition-garbage failure signature once forced past the same threshold
on a single, ordinary generation call. Since HotpotQA's per-hop contexts
(1{,}500--2{,}500+ tokens) routinely push a three-agent relay chain's
cumulative cache past this threshold, this rules Phi-3.5-mini-instruct
out for this pipeline regardless of further pipeline-side fixes, rather
than indicating a bug we could still close before submission. Whether
findings -- particularly layer-selection's architecture-dependence --
extend to substantially larger or smaller models, mixture-of-experts
architectures, or models outside this decoder-only,
grouped-query-attention, RoPE lineage remains largely untested beyond
this one additional data point (one further methodological risk worth
flagging: \citet{latentcommpay2026} reports at least one model pairing
where the receiver could not consume a relayed cache at all, i.e.
cross-architecture relay can fail at the precondition stage, not just on
accuracy -- a null result under those conditions would itself be
informative, not a bug). All of our own architecture variation is
\emph{between} pipeline runs (the same model family fills every agent
role within a given run); whether a single pipeline mixing architectures
across agent roles -- a Qwen Reasoner handing its cache to a Mistral
Verifier, for instance -- is even viable is untested and, to our
knowledge, an open problem in this literature more broadly: KVCOMM itself
flags relay between ``agents with identical architectures but different
weights'' and agents with ``different attention formulations'' as work it
leaves for future exploration, and we inherit that same gap rather than
closing it. Nearly all results are on HotpotQA
(distractor configuration); we ran one $n=50$ confirmatory check on GSM8K
(Qwen, \texttt{A}/\texttt{D}/full causal audit on \texttt{A}) and both
headline findings replicated cleanly -- \texttt{D} matched \texttt{A}'s
accuracy exactly (McNemar $p=1.0$, more decisively free than on
HotpotQA) and the causal audit's three-tier ordering held with every
pairwise comparison significant ($p<0.0002$), larger effect sizes than the
original HotpotQA run. This is one run on one additional task family, not
a systematic multi-task study -- we did not extend it to
\texttt{B\_int8}/\texttt{D}-family causal audits or Mistral, and report it
as suggestive, not as evidence the claims generalize broadly across
tasks. This matters specifically for the offset-correction negative
result: KVCOMM was evaluated on MMLU, GSM8K, and HumanEval, none of which
is multi-hop QA, and whether the technique fares differently on the more
homogeneous task family it was originally validated on remains open (our
own GSM8K check did not re-test \texttt{E}). All results additionally use
greedy decoding (\texttt{do\_sample=False}); we did not evaluate under
sampling or self-consistency, which could interact differently with a
degraded or compressed relayed context than with a full one.

\paragraph{Nibble-packing is unit-tested, not deployment-tested.} We
implement real two-values-per-byte packing for every bits==4 tensor (see
Finding 4), verified via round-trip pack/unpack exactness and
bit-identical \texttt{compress}/\texttt{decompress} output with and
without packing on synthetic tensors. This is standalone tensor-level
testing, not an end-to-end GPU rerun with the new packing path in the
live pipeline; accuracy/F1 are unaffected in principle (packing changes
storage format, not values, and the unit tests confirm this directly),
but we have not re-run the full $n=100$ HotpotQA evaluation with packing
enabled to confirm no integration issue exists between the compressor and
the rest of the pipeline (cache reconstruction, device placement, batch
handling). The reported accuracy numbers throughout this paper come from
pre-packing runs; only the compression-ratio figures reflect the
packing-enabled compressor.

\paragraph{Partial reproduction of adapted and considered techniques.}
Our rotation-based quantization covers PolarQuant's core
rotate-then-quantize mechanism but omits TurboQuant's QJL bias-correction
term; our KIVI-inspired fix (Section~\ref{sec:method}) omits the original
method's group-wise windowing and fp16 residual buffer, implementing
full-sequence per-channel quantization only -- the minimal change
isolating the RoPE-interaction mechanism we diagnosed, not a complete
reproduction. We report both as inspired-by implementations scoped to
test a specific hypothesis, not certified reproductions, and encourage
readers wanting the original techniques' full guarantees to consult the
primary sources. We considered but did not implement a fourth
layer-reconstruction strategy, Orthogonal Backfill
\citep{orthobackfill2026}: its published formula requires attention
weights this pipeline's \texttt{sdpa} decode path does not expose, and
given our own evidence that naive/simplified reproductions can actively
mislead (the rotation-based quantization finding above), we chose not to
ship an approximated version under time pressure -- a scoping decision,
not a result.

We checked ARR's guidance before deciding whether a separate Ethical
considerations section was warranted: this work uses two long-established
public benchmarks (HotpotQA, GSM8K), two openly-released model
checkpoints, no human subjects, no crowdsourcing, and no new data
collection -- we are not aware of a concern specific to this paper beyond
what this section already covers.

\bibliography{references_rlm_focus}
\bibliographystyle{acl_natbib}

\end{document}
